\documentclass[11pt,a4paper]{article}
% Shared preamble for RuPhO English problem statements (match T1 2026 style).
% Use: \input{latex/rupho_en_preamble.tex} after \documentclass.
\usepackage[margin=1in]{geometry}
\usepackage{amsmath,amssymb}
\usepackage{graphicx}
\usepackage{float}
\usepackage{enumitem}
\usepackage{microtype}
\usepackage{caption}
\usepackage{tabularx}
\usepackage{hyperref}

\captionsetup{font=small,labelfont=bf,skip=6pt}

\setlist[itemize]{leftmargin=1.2cm}
\setlist[enumerate]{leftmargin=1.2cm}

% One outer box per subproblem: [id | statement | points] (no inner rules)
\newcommand{\problembox}[3]{%
  \par\addvspace{0.42em}%
  \noindent
  \setlength{\fboxsep}{7.5pt}%
  \setlength{\fboxrule}{0.95pt}%
  \fbox{%
    \begin{minipage}{\dimexpr\linewidth-2\fboxsep-2\fboxrule\relax}
    \begin{tabularx}{\linewidth}{@{}>{\raggedright\arraybackslash\bfseries}p{2.1em} @{\hspace{0.9em}} >{\raggedright\arraybackslash}X @{\hspace{0.9em}} >{\raggedleft\arraybackslash\bfseries}p{2.4em}@{}}
    #1 & #2 & #3
    \end{tabularx}
    \end{minipage}%
  }%
  \par\addvspace{0.32em}%
}

\title{\textbf{''Uncomplicated'' capacitor}}
\author{\normalsize W22 (2022) --- English translation}
\date{}
\begin{document}
\maketitle

\section*{Part A. Charged plate in a capacitor.}

In the middle between the identical uncharged metal plates of a flat capacitor with an area $S$ and a distance between them $d\ll{\sqrt{S}}$ there is a dielectric plate of mass $m$ of the same shape, charged uniformly over the surface with a charge $q>0$. The capacitor plates were connected to a constant voltage source $\mathcal{E}$ with zero internal resistance. Then the plate is released. When moving, the capacitor plates remain motionless. Let us denote by $x$ the deviation of the plate from the initial position, by $q_1$ the charge of the left plate, and by $q_2$ the charge of the right plate.

\begin{figure}[H]
    \centering
    \includegraphics[width=0.72\textwidth]{T1_EN_assets/fig1.jpg}
    \caption{Figure 1.}
\end{figure}

\problembox{A1}{Express the charges of the capacitor plates $q_1(x)$ and $q_2(x)$ in terms of $\mathcal{E}$, $S$, $d$, $q$, $x$ and physical constants.}{0.50}

\problembox{A2}{Express the force of electrostatic interaction $F(x)$ of the plate with the plates in terms of $\mathcal{E}$, $S$, $d$, $q$, $x$ and physical constants.}{0.50}

It can be shown that the dependence of the displacement of the plate $x(t)$ on time has the following form $$x(t)=Ae^{at}+Be^{-at}+C, $$ where $a$   from $\mathcal{E}$, $S$, $d$, $q$ and $m$,  $A$, $B$ and $C$ are determined by the initial conditions.

\problembox{A3}{Express $a$, $A$, $B$ and $C$ in terms of $\mathcal{E}$, $S$, $d$, $q$, $m$ and physical constants.}{1.00}

\problembox{A4}{Express the speed of the plate $v^*$ at the moment of impact on the capacitor plate in terms of $\mathcal{E}$, $S$, $d$, $q$, $m$ and physical constants.}{1.00}

\section*{Part B. Capacitor with Dielectric}

A flat capacitor with plate area $S$ and distance between them $d\ll{\sqrt{S}}$ is filled with a dielectric with dielectric constant $\varepsilon$ and resistivity $\rho$. The dielectric and the plates have good electrical contact. Let the potential difference between the plates of the capacitor be $V$.

\begin{figure}[H]
    \centering
    \includegraphics[width=0.72\textwidth]{T1_EN_assets/fig2.jpg}
    \caption{Figure 2.}
\end{figure}

\problembox{B1}{Find the charges $q$ on the metal plates of the capacitor. Express your answer in terms of $V$, $S$, $d$, $\varepsilon$ and physical constants.}{0.30}

\problembox{B2}{Find the current $I$ flowing through the dielectric. Express your answer in terms of $V$, $S$, $d$, $\rho$ and physical constants.}{0.30}

From the two previous points it follows that the circuit of such a capacitor in an electrical circuit is equivalent to some connection of a resistor with resistance $R$ and a capacitor with capacitance $C$.

\problembox{B3}{Find $R$ and $C$, and also indicate the type of connection of these elements (series or parallel).}{0.40}

Now the complex capacitor is connected to a constant voltage source $\mathcal{E}$, which has an internal resistance $r$. At the initial moment of time, the capacitor is not charged.

\problembox{B4}{Find the time dependence of the current $I(t)$ flowing through the source. Express your answer in terms of $\mathcal{E}$, $r$, $R$, $C$ and $t$.}{2.00}

\section*{Part C. Two-Dielectric Capacitor}

A flat capacitor with a distance between plates $d$ is filled with two dielectrics. One with dielectric constant $\varepsilon_1$, resistivity $\rho_1$ and contact area with plates $S_1$, and the other with dielectric constant $\varepsilon_2$, resistivity $\rho_2$ and contact area with plates $S_2$. Both dielectrics have good electrical contact with the capacitor plates.

\problembox{C1}{Find the equivalent capacitance $C_0$ of this capacitor. Express your answer in terms of $S_1$, $S_2$, $\varepsilon_1$, $\varepsilon_2$, $d$ and physical constants.}{0.40}

\problembox{C2}{Find the equivalent resistance $R_0$ of this capacitor. Express your answer in terms of $S_1$, $S_2$, $\rho_1$, $\rho_2$ and $d$.}{0.40}

This capacitor in series with the capacitor $C$ is connected to a constant voltage source $\mathcal{E}$, which has zero internal resistance. While the key was open, the capacitors were not charged. The voltage across a complex capacitor is measured with an ideal voltmeter $V$.

\begin{figure}[H]
    \centering
    \includegraphics[width=0.72\textwidth]{T1_EN_assets/fig3.jpg}
    \caption{Figure 3.}
\end{figure}

\problembox{C3}{Find the charge $q_C(0)$ on the capacitor capacitance $C$ immediately after closing the switch. Express your answer using $\mathcal{E}$, $C$ and $C_0$.}{0.40}

\problembox{C4}{Find the dependence of the voltmeter readings $V(t)$ on time $t$ after closing the key. Express your answer in terms of $\mathcal{E}$, $C$, $C_0$, $R_0$ and $t$.}{1.30}

\problembox{C5}{How much heat $Q_{R_0}$ will be released at resistance $R_0$? Express your answer using $\mathcal{E}$, $C$ and $C_0$.}{0.50}

$\textit{}$ : $$\textbackslash{}int\textbackslash{}frac{dx}{ax+b}=\textbackslash{}frac{1}{a}\textbackslash{}ln\textbackslash{}left(ax+b\textbackslash{}right)+C $$ $$\textbackslash{}int e^{ax}dx=\textbackslash{}frac{e^{ax}}{a}+C $$

\vspace{1em}
\noindent\textit{Source:} \url{https://pho.rs/p/271}

\end{document}